\documentclass[11pt]{article}
\usepackage[margin=1in]{geometry}
\usepackage{amsmath,amssymb,amsthm,hyperref}
\newtheorem{theorem}{Theorem}
\newtheorem{lemma}{Lemma}
\title{Logarithmic biorthogonal norms for the Gaussian Gabor system\\
\large MathDB 375037}
\author{Principia Math in Progress}
\date{12 September 2026}
\begin{document}
\maketitle

\paragraph{Task 1.}
Prove Conjecture 4.1 of Zikkos \cite{Z}: after deleting the Gaussian Gabor
atom indexed by $(0,0)$, the unique biorthogonal vectors satisfy
\begin{equation}\label{eq:target}
 \|\gamma_{m,n,0,0}\|_{L^2(\mathbb R)}^2
 \asymp\log\!\left(e+\sqrt{m^2+n^2}\right),
 \qquad (m,n)\in\mathbb Z^2\setminus\{(0,0)\}.
\end{equation}
Also prove the stated consequence for positive weights whose
logarithmic squared-reciprocal sum is finite.
All comparison constants below are independent of $(m,n)$.

\begin{theorem}
Estimate \eqref{eq:target} holds. If $c_{m,n}>0$ and
\begin{equation}\label{eq:weights}
 \sum_{(m,n)\ne(0,0)}
 \frac{\log(e+\sqrt{m^2+n^2})}{c_{m,n}^2}<\infty,
\end{equation}
then the family
\[
 \left\{c_{m,n}e^{2\pi int}e^{-\pi(t-m)^2}:
              (m,n)\in\mathbb Z^2\setminus\{(0,0)\}\right\}
\]
is a complete Riesz--Fischer sequence and a lower semi-frame for
$L^2(\mathbb R)$.
\end{theorem}

\paragraph{The Zak transform and its zero.}
Initially normalize the Gaussian as $\phi(t)=2^{1/4}e^{-\pi t^2}$.
The Zak transform
\[
 Zf(x,\xi)=\sum_{j\in\mathbb Z}f(x-j)e^{2\pi ij\xi}
\]
is unitary from $L^2(\mathbb R)$ onto $L^2(Q)$, where $Q=[0,1)^2$.
Write $\Theta=Z\phi$ and
$E_{m,n}(x,\xi)=e^{2\pi i(nx-m\xi)}$. Changing the summation index gives
\begin{equation}\label{eq:covariance}
 Z\big(e^{2\pi int}\phi(t-m)\big)=E_{m,n}\Theta.
\end{equation}

The classical theta-function description gives a unique zero of $\Theta$
in $Q$, at $q_0=(1/2,1/2)$. It is a nondegenerate zero:
\begin{equation}\label{eq:zero}
 c_0|u|\leq|\Theta(q_0+u)|\leq C_0|u|
 \quad (|u|<\delta)
\end{equation}
for positive $c_0,C_0$ and $0<\delta<1/2$.
For clarity, the published proof in \cite[Corollary 4.3, equation (11)]{HHKN}
gives the expansion
\[
 \Theta(q_0+u)
 =-2^{1/4}\pi\theta_1'(0,e^{-\pi})(u_1+iu_2)+O(|u|^2),
 \qquad \theta_1'(0,e^{-\pi})\ne0.
\]
The derivative is nonzero also by the Jacobi triple product.
The modulus of $\Theta$ is periodic in both coordinates and bounded away
from zero on the complement of this disk in the closed unit square.
In particular, $1/\Theta\notin L^2(Q)$.

\paragraph{The biorthogonal formula.}
For $(m,n)\ne(0,0)$ let
\[
 a_{m,n}=E_{m,n}(q_0)=(-1)^{n-m},\qquad
 H_{m,n}=\frac{E_{m,n}-a_{m,n}}{\overline\Theta}.
\]
The numerator vanishes to at least first order at $q_0$. By
\eqref{eq:zero}, this quotient is locally bounded for each fixed index,
and it is bounded away from the zero as well. Thus $H_{m,n}\in L^2(Q)$.
Let $h_{m,n}=Z^{-1}H_{m,n}$. For all nonzero $(r,s)$, Fourier
orthogonality gives
\begin{align}
 \left\langle e^{2\pi ist}\phi(t-r),h_{m,n}\right\rangle
 &=\int_Q E_{r,s}\Theta\,
       \overline{\frac{E_{m,n}-a_{m,n}}{\overline\Theta}}\,dA\notag\\
 &=\int_Q E_{r,s}\overline{E_{m,n}}\,dA
   -\overline{a_{m,n}}\int_Q E_{r,s}\,dA
 =\delta_{r,m}\delta_{s,n}.                    \label{eq:biorth}
\end{align}
The constant term disappears because the origin index was deleted.

For completeness, if $F\in L^2(Q)$ is orthogonal to all
$E_{r,s}\Theta$ with $(r,s)\ne(0,0)$, the function $F\overline\Theta$
has only a constant Fourier coefficient. It is therefore constant
almost everywhere. Since $1/\Theta\notin L^2(Q)$, the constant must be
zero, and $F=0$ almost everywhere. Thus the deleted system is complete.
The exhibited biorthogonal vectors imply minimality, and completeness
implies uniqueness of the biorthogonal family.
This recovers the known formula in \cite[Theorem 4.2]{HHKN}. In particular,
\begin{equation}\label{eq:norm}
 \|h_{m,n}\|_2^2
 =\int_Q\frac{|E_{m,n}-a_{m,n}|^2}{|\Theta|^2}\,dA.
\end{equation}

\begin{lemma}\label{lem:integral}
For fixed $\delta>0$, define
\[
 J_\delta(p)=\int_{|u|<\delta}
       \frac{|e^{2\pi ip\cdot u}-1|^2}{|u|^2}\,dA(u).
\]
As $|p|\to\infty$, $J_\delta(p)\asymp\log|p|$, uniformly in the
direction of $p\in\mathbb R^2$.
\end{lemma}
\begin{proof}
Rotate polar coordinates so that $p$ points along the first axis and
put $R=|p|$. With
$J_0(s)=(2\pi)^{-1}\int_0^{2\pi}\cos(s\cos\theta)\,d\theta$,
the angular integral gives
\begin{equation}\label{eq:bessel}
 J_\delta(p)=4\pi\int_0^{2\pi\delta R}\frac{1-J_0(s)}s\,ds.
\end{equation}
The inequality $1-\cos t\leq t^2/2$ implies
$0\leq1-J_0(s)\leq s^2/4$ near zero.
For every $s>0$, $J_0(s)<1$, since the cosine in its angular average
is not identically one. Moreover,
\[
 J_0(s)=\frac1\pi\int_{-1}^{1}\frac{\cos(st)}{\sqrt{1-t^2}}\,dt
 \longrightarrow0
\]
by the Riemann--Lebesgue lemma. Continuity and compactness therefore
give a fixed $\eta>0$ such that
$\eta\leq1-J_0(s)\leq2$ for every $s\geq1$.
The integral over $(0,1)$ in \eqref{eq:bessel} is bounded, while that
over $(1,2\pi\delta R)$ is comparable to $\log R$.
The formula depends only on $R$, giving uniformity in direction.
\end{proof}

\paragraph{Completion of the estimate.}
In the disk around $q_0$,
\[
 E_{m,n}(q_0+u)=a_{m,n}e^{2\pi i(nu_1-mu_2)}.
\]
Equations \eqref{eq:zero} and \eqref{eq:norm} sandwich the local norm
contribution between fixed positive multiples of
$J_\delta((n,-m))$. On the complement, the numerator is at most two in
modulus and $|\Theta|$ is bounded below, so the remaining contribution
is bounded uniformly in $(m,n)$.
Lemma \ref{lem:integral} gives
$\|h_{m,n}\|_2^2\asymp\log\sqrt{m^2+n^2}$ for large radius.
There are only finitely many nonzero lattice points in a bounded disk,
and each corresponding norm is positive. Adjusting the constants
therefore proves the uniform bound with $\log(e+\sqrt{m^2+n^2})$.

The source uses $e^{-\pi t^2}=2^{-1/4}\phi(t)$, so its biorthogonal
vectors are $\gamma_{m,n,0,0}=2^{1/4}h_{m,n}$.
Their squared norms differ by the fixed factor $\sqrt2$.
This proves \eqref{eq:target}.

Finally, \eqref{eq:weights} implies
$\sum\|\gamma_{m,n,0,0}\|_2^2/c_{m,n}^2<\infty$.
The deleted system is complete and minimal, as proved above.
Theorem 1.2 of \cite{Z} now gives both the Riesz--Fischer and lower
semi-frame assertions, including completeness.
\hfill$\square$

\paragraph{Prior results and verification.}
The explicit biorthogonal formula and the simple Zak zero are known.
Belov \cite[equation (3.1)]{B} already proved the logarithmic upper
bound in Fock coordinates; \cite[Lemma 2.3]{BBB} repeats it.
The argument here supplies the uniform matching lower bound and the
complete two-sided estimate requested in the 2026 conjecture.
Exact Fourier checks verify the index and phase conventions; supplementary
numerical checks compare the Gaussian Zak series with the theta formula
and independently check the angular integral. Finite computations are
not used to infer the asymptotic statement.

\paragraph{Task 2.}
Determine the leading constant in the logarithmic norm asymptotic and
prove a uniform bounded remainder using the local Taylor expansion
of the Gaussian Zak transform.

\paragraph{Task 3.}
Find necessary and sufficient growth conditions for radial weights
$c_{m,n}=c(\sqrt{m^2+n^2})$ to produce a lower semi-frame, beyond the
sufficient summability condition \eqref{eq:weights}.

\paragraph{Task 4.}
Extend the norm-growth analysis to other windows whose Zak transforms
have isolated zeros, including higher-order zeros and the deletion
of an appropriate finite set of atoms.

\paragraph{Task 5.}
Obtain quantitative stability bounds for finite reconstruction from the
weighted Gaussian system, relating the truncation radius and weight
growth to the norms of the biorthogonal vectors.

\begin{thebibliography}{9}
\bibitem{Z}
E. Zikkos, \emph{The Gaussian Gabor system at the critical density is a
weighted lower semi frame}, 2026.
\href{https://arxiv.org/abs/2606.00764}{arXiv:2606.00764}, Conjecture 4.1.
\bibitem{HHKN}
L. Hart, C. Heil, I. Katz, and M. Northington V,
\emph{Overcomplete reproducing pairs},
J. Fourier Anal. Appl. \textbf{31} (2025), article 58.
\href{https://heil.math.gatech.edu/papers/reproducingpairs_journal.pdf}{Published text}.
\bibitem{B}
Y. Belov, \emph{Uniqueness of Gabor series},
Appl. Comput. Harmon. Anal. \textbf{39} (2015), 545--551.
\href{https://arxiv.org/abs/1409.5231}{arXiv:1409.5231}.
\bibitem{BBB}
A. Baranov, Y. Belov, and A. Borichev,
\emph{Summability properties of Gabor expansions},
J. Funct. Anal. \textbf{274} (2018), 2532--2552.
\href{https://arxiv.org/abs/1706.05685}{arXiv:1706.05685}.
\end{thebibliography}
\end{document}
